\clearpage\pagestyle{otMain}\setlength{\otBlocEspacement}{2.5mm}\fontsize{11}{15}\selectfont\corrigesChapitre{20}{Organiser et lire des données}
\begin{blocLivre}{cor-20-1}\begin{corrige}{1}{Lire un tableau}\demarche{Lire la bonne colonne puis additionner.}
La colonne fraise indique $\res9$ élèves. Le total vaut $7+9+5+4=\res{25}$ élèves.
\end{corrige}\end{blocLivre}
\begin{blocLivre}{cor-20-2}\begin{corrige}{2}{Double entrée}\demarche{Croiser ligne et colonne avant d’additionner.}
Au croisement de la ligne CM2 et de la colonne Menu A : $\res{14}$ élèves. Menu B : $8+12=\res{20}$ élèves.
\end{corrige}\end{blocLivre}
\begin{blocLivre}{cor-20-3}\begin{corrige}{3}{Compléter les totaux}\demarche{Vérifier le total par lignes puis par colonnes.}
CM1 : $17+8=\res{25}$. CM2 : $14+12=\res{26}$. Total : $25+26=\res{51}$ élèves ; on retrouve ce total en additionnant les colonnes : $31+20=51$.
\end{corrige}\end{blocLivre}
\begin{blocLivre}{cor-20-4}\begin{corrige}{4}{Le diagramme en barres}\demarche{Lire la hauteur de chaque barre sur l’axe.}
Mardi : $\res{12}$ livres ; jeudi : $\res{18}$ livres, la barre la plus haute : c’est le $\res{\text{jeudi}}$. Chaque graduation vaut 2 livres.
\label{fig-cor-20-4}\figureGeo[corrige]{don-4}
\end{corrige}\end{blocLivre}
\begin{blocLivre}{cor-20-5}\begin{corrige}{5}{La courbe de température}\demarche{Suivre la verticale de l’heure jusqu’à la courbe.}
À 8 h : $\res{12\,^\circ\mathrm C}$ ; à 14 h : $\res{21\,^\circ\mathrm C}$, le point le plus haut de la courbe : le maximum est à $\res{14\ \text h}$.
\label{fig-cor-20-5}\figureGeo[corrige]{don-5}
\end{corrige}\end{blocLivre}
\begin{blocLivre}{cor-20-6}\begin{corrige}{6}{Construire un diagramme}\demarche{Une barre par catégorie, hauteur égale à l’effectif.}
Chaque fruit reçoit une barre de la hauteur de son effectif : 7, 9, 5 et 4. La barre la plus haute est celle de la fraise. Le diagramme figure en vert.
\label{fig-cor-20-6}\figureGeo[corrige]{don-6}
\end{corrige}\end{blocLivre}
\begin{blocLivre}{cor-20-7}\begin{corrige}{7}{Des points dans un repère}\demarche{Lire les deux coordonnées dans le bon ordre.}
À 2 semaines : $\res{2\ \mathrm{kg}}$. À 6 semaines : $\res{5\ \mathrm{kg}}$. On lit l’abscisse (l’âge) puis l’ordonnée (la masse) de chaque point.
\label{fig-cor-20-7}\figureGeo[corrige]{don-7}
\end{corrige}\end{blocLivre}
\begin{blocLivre}{cor-20-8}\begin{corrige}{8}{Des questions en étapes}\demarche{Extraire les bonnes cases avant de calculer.}
Menu B : $12-8=\res4$ élèves de plus en CM2. Menu A : $17+14=31$ choix contre 20 : $\res{\text{oui}}$, avec $31-20=11$ choix d’écart.
\end{corrige}\end{blocLivre}
\begin{blocLivre}{cor-20-9}\begin{corrige}{9}{L’année des températures}\demarche{Chercher les points extrêmes de la courbe.}
Le point le plus haut est en $\res{\text{juillet}}$ ($22\,^\circ$C), le plus bas en $\res{\text{janvier}}$ ($4\,^\circ$C). L’écart vaut $22-4=\res{18}$ degrés.
\label{fig-cor-20-9}\figureGeo[corrige]{don-9}
\end{corrige}\end{blocLivre}
\begin{blocLivre}{cor-20-10}\begin{corrige}{10}{La barre fausse}\demarche{Comparer chaque barre à sa valeur du tableau.}
La barre du $\res{\text{mardi}}$ monte à 10 au lieu de $\res{12}$. Les barres de lundi et mercredi sont correctes ; la barre corrigée figure en vert.
\label{fig-cor-20-10}\figureGeo[corrige]{don-10}
\end{corrige}\end{blocLivre}
\begin{blocLivre}{cor-20-11}\begin{corrige}{11}{L’enquête brute}\demarche{Barrer les réponses au fur et à mesure du comptage.}
En comptant : à pied $\res5$, vélo $\res3$, bus $\res2$. Le total $5+3+2=10$ retrouve les dix réponses. Le moyen le plus utilisé est la marche $\res{\text{à pied}}$.
\end{corrige}\end{blocLivre}
\begin{blocLivre}{cor-20-12}\begin{corrige}{12}{La bibliothèque}\demarche{Choisir les cases utiles à chaque question.}
Semaine 1 : $34+21=\res{55}$ prêts. Documentaires : de 21 à 25, ils ont $\res{\text{augmenté}}$ de 4. Romans : $34+28=\res{62}$ en deux semaines.
\end{corrige}\end{blocLivre}
\begin{blocLivre}{cor-20-13}\begin{corrige}{13}{La plante qui pousse}\demarche{Suivre la ligne horizontale des 10 cm.}
À la semaine 3 : $\res{6\ \mathrm{cm}}$. La courbe franchit la ligne des 10 cm entre les semaines 4 et 5 : la plante dépasse 10 cm à partir de la $\res{\text{semaine }5}$, où elle atteint 13 cm.
\label{fig-cor-20-13}\figureGeo[corrige]{don-13}
\end{corrige}\end{blocLivre}
\begin{blocLivre}{cor-20-14}\begin{corrige}{14}{Deux villes, deux courbes}\demarche{Comparer les courbes saison par saison.}
En été, la courbe de Lyonnet est au-dessus : $\res{\text{Lyonnet}}$ ($24\,^\circ$C contre 18). En hiver, Brestel garde $\res{7\,^\circ\mathrm C}$ contre 3. Les variations de Brestel, de 7 à 18 degrés, sont plus resserrées : ses températures sont $\res{\text{les plus régulières}}$.
\label{fig-cor-20-14}\figureGeo[corrige]{don-14}
\end{corrige}\end{blocLivre}
\begin{blocLivre}{cor-20-15}\begin{corrige}{15}{Les précipitations}\demarche{Répondre à chaque question avec les bonnes cases.}
Le mois le plus pluvieux est $\res{\text{novembre}}$ (96 mm). Total : $45+78+96+60=\res{279}$ mm. Septembre et décembre réunis : $45+60=105$ mm ; novembre, avec 96 mm, en a reçu $\res{\text{moins}}$.
\end{corrige}\end{blocLivre}
